\documentclass[11pt, a4paper]{report}

\usepackage[margin=2.5cm]{geometry}
\usepackage[T1]{fontenc}
\usepackage[utf8]{inputenc}
\usepackage{amsmath, amssymb, amsthm}
\usepackage{graphicx}
\usepackage{subcaption}
\usepackage{booktabs}
\usepackage{multirow}
\usepackage{float}
\usepackage{listings}
\usepackage{xcolor}
\usepackage{enumitem}
\usepackage{hyperref}
\usepackage{url}
\usepackage{caption}
\usepackage{setspace}
\usepackage{csquotes}
\usepackage{siunitx}

\definecolor{codebg}{RGB}{245,245,245}
\definecolor{commentgray}{RGB}{90,90,90}
\definecolor{keywordblue}{RGB}{30,76,145}
\definecolor{stringgreen}{RGB}{20,100,40}

\lstset{
    backgroundcolor=\color{codebg},
    basicstyle=\ttfamily\small,
    keywordstyle=\color{keywordblue}\bfseries,
    commentstyle=\color{commentgray}\itshape,
    stringstyle=\color{stringgreen},
    breaklines=true,
    frame=single,
    numbers=left,
    numberstyle=\tiny,
    showstringspaces=false,
    tabsize=2,
    columns=fullflexible
}

\onehalfspacing

\title{\textbf{Bird Species Classification using CNN A and CNN B}\\[0.5em] \Large Preliminary Technical Report}
\author{Your Name}
\date{\today}

\begin{document}

\maketitle
\tableofcontents
\listoffigures
\listoftables

\chapter{Introduction}
\section{Project Aim}
This report summarises the development and evaluation of two convolutional neural network pipelines for bird image classification. The first model, \textbf{CNN A}, was developed as a task-specific architecture trained from scratch on the reduced-resolution dataset. The second model, \textbf{CNN B}, was created using transfer learning from the best-performing CNN A model and then adapted to the full-resolution image task.

\begin{itemize}[leftmargin=1.5em]
    \item Briefly state the classification problem and the dataset context.
    \item Explain the motivation for comparing a lightweight custom CNN with a transferred model.
    \item State the research objective: identify the most suitable architecture for the task while balancing performance and computational cost.
\end{itemize}

\section{Summary of Methods}
\begin{itemize}[leftmargin=1.5em]
    \item Part 1: train and tune a compact CNN using reduced-resolution images to identify promising architectures.
    \item Select the best-performing configuration as \textbf{CNN A}.
    \item Retrain and evaluate the strongest variants on the full-resolution dataset.
    \item Part 2: load the best model from Part 1 and adapt it via transfer learning to form \textbf{CNN B}.
    \item Compare CNN A and CNN B using validation and test performance, along with estimated computational effort.
\end{itemize}

\chapter{Dataset and Preprocessing}
\section{Data Description}
\begin{itemize}[leftmargin=1.5em]
    \item Mention the number of classes and the bird image dataset structure.
    \item State that data are split into train, validation and test sets.
    \item Describe the use of class folders and the class-label mapping produced by \texttt{flow\_from\_directory}.
\end{itemize}

\section{Image Preprocessing}
\begin{itemize}[leftmargin=1.5em]
    \item All images were rescaled to $[0,1]$ using $1/255$.
    \item Input sizes were set to $128 \times 128$ for reduced-resolution searches and $224 \times 224$ for final evaluation.
    \item Data augmentation was applied during training where appropriate (rotation, shifts, zoom, shear, horizontal flip).
    \item For class imbalance, class weights were considered and, in some cases, used instead of oversampling.
\end{itemize}

\section{Class Imbalance}
The training split showed mild imbalance across species. This is relevant because some bird classes are represented more frequently than others. To avoid artificially inflating the minority classes by duplicating samples, a weighted loss formulation was used when necessary.

\begin{equation}
\mathcal{L} = -\sum_{i=1}^{C} w_i \, y_i \log(\hat{y}_i)
\end{equation}

where $C$ is the number of classes, $y_i$ is the ground-truth label, $\hat{y}_i$ is the predicted probability, and $w_i$ is the class weight for class $i$.

\chapter{Model Architectures}
\section{CNN A: Best Model from Part 1}
CNN A was selected as the strongest candidate from the reduced-resolution hyperparameter search and then retrained on the full-resolution data. Its architectural pattern was designed to provide a strong balance between feature extraction and efficiency.

\begin{itemize}[leftmargin=1.5em]
    \item Use of repeated convolutional blocks with \texttt{Conv2D} followed by \texttt{BatchNormalization} and \texttt{MaxPooling2D}.
    \item Moderate number of filters to retain discriminative features without excessive memory cost.
    \item Final pooling layer followed by a dense classifier and a softmax output.
    \item Learning rate and dropout tuned using the validation signal.
\end{itemize}

The architecture was chosen to satisfy three goals:
\begin{enumerate}[leftmargin=2em]
    \item strong feature extraction for bird morphology,
    \item manageable computational cost,
    \item stable convergence under the limited tuning budget.
\end{enumerate}

\section{CNN B: Transfer-Learned Model from Part 2}
CNN B was built from the best-performing CNN A model by freezing the base feature extractor and attaching a new task-specific classification head. This approach is useful when the feature space learned on the source task is transferable to the bird classification task.

\begin{itemize}[leftmargin=1.5em]
    \item Base model loaded from the saved CNN A checkpoint.
    \item Early layers kept frozen to retain general visual features.
    \item A new dense head was trained on the bird dataset.
    \item Optional fine-tuning of a later convolution block was explored if validation performance plateaued.
\end{itemize}

\chapter{Data Augmentation and Training Strategy}
\section{Augmentation Strategy}
\begin{itemize}[leftmargin=1.5em]
    \item Rotation, width/height shifts, shear, zoom, and horizontal flips were used for training augmentation.
    \item Augmentation was used primarily to improve robustness to pose and viewpoint variation.
    \item The augmentation policy was not treated as a substitute for class weighting, since the main problem was mild class imbalance rather than lack of diversity alone.
\end{itemize}

\section{Training Procedure}
\begin{itemize}[leftmargin=1.5em]
    \item Hyperparameter tuning was performed with a reduced-resolution search to identify promising configurations efficiently.
    \item The best models were retrained at full image resolution for final comparison.
    \item Early stopping was used to prevent unnecessary training and reduce overfitting.
    \item Validation accuracy and validation loss were used to select the best checkpoint.
\end{itemize}

\chapter{Results and Findings}
\section{Comparative Performance}
The final comparison focused on whether CNN B improved over CNN A in terms of predictive quality and training efficiency.

\begin{itemize}[leftmargin=1.5em]
    \item CNN A is expected to be strong in terms of task-specific optimisation and relatively low training cost.
    \item CNN B can benefit from feature reuse and may converge faster in early epochs.
    \item CNN B may outperform CNN A on a held-out test set if the transferred features transfer well to the bird domain.
    \item CNN A may still be competitive because it was tailored directly to the final task and may require fewer optimisation steps after selection.
\end{itemize}

\section{Performance vs Computational Effort}
This comparison should clearly distinguish between two types of efficiency:
\begin{itemize}[leftmargin=1.5em]
    \item \textbf{Performance}: classification accuracy, validation accuracy, and test accuracy.
    \item \textbf{Computational effort}: training time, GPU memory, and number of epochs required for convergence.
\end{itemize}

A model can be better in accuracy but worse in computational cost if it requires longer training or larger input sizes. In this project, the central question is whether CNN B provides a material gain in accuracy sufficient to justify the additional transfer-learning training overhead.

\section{Expected Outcome}
\begin{itemize}[leftmargin=1.5em]
    \item CNN A is likely to be a strong and efficient custom model for the task.
    \item CNN B is likely to have a practical advantage when the base model already captures useful visual features.
    \item The final conclusion should be based on the observed validation/test performance and the training time required to reach that performance.
\end{itemize}

\chapter{Resource Estimate for the Full 525-Species Task}
For the full bird classification problem with 525 species, the resource requirement would increase substantially compared with the current experiment.

\begin{equation}
\text{Estimated compute} \propto N_{samples} \times C \times H \times W \times E \times B
\end{equation}

where:
\begin{itemize}[leftmargin=1.5em]
    \item $N_{samples}$ is the number of training images,
    \item $C$ is the number of classes,
    \item $H \times W$ is the image resolution,
    \item $E$ is the number of epochs,
    \item $B$ is the batch size or effective training workload.
\end{itemize}

\begin{itemize}[leftmargin=1.5em]
    \item The class count alone rises from 10 to 525, increasing the output space and the likelihood of imbalance.
    \item A larger number of classes requires stronger regularisation and more careful class balancing.
    \item Full-resolution training on 525 classes will likely require more GPU memory and longer training time.
    \item Transfer learning is strongly recommended for the full task, because it reduces the need to train a deep CNN from scratch.
    \item A realistic scale-up plan would include a pretrained backbone, class-weighting, and a longer validation-driven training schedule.
\end{itemize}

\chapter{Discussion}
\section{Why CNN A and CNN B Were Chosen}
\begin{itemize}[leftmargin=1.5em]
    \item CNN A was chosen because it performed best under the reduced-resolution tuning and provided a compact but expressive model family.
    \item CNN B was chosen because transfer learning reuses strong visual features and is often more efficient than training a deeper network from random initialisation.
    \item The comparison is useful because it distinguishes between a task-specific handcrafted architecture and a pretrained feature extractor adapted to the final task.
\end{itemize}

\section{Interpretation of Results}
\begin{itemize}[leftmargin=1.5em]
    \item If CNN B achieves better test accuracy with similar or lower training time, it should be treated as the preferred final model.
    \item If CNN A is nearly as accurate but cheaper to train, it may be preferable for deployment in resource-constrained settings.
    \item A model that is computationally cheaper but slightly less accurate may still be attractive for iterative experimentation and rapid retraining.
\end{itemize}

\chapter{Conclusion}
This project compared a custom CNN trained from scratch against a transfer-learned CNN based on the best configuration from Part 1. The two models are complementary: CNN A provides a task-specific baseline, while CNN B tests whether pretrained feature reuse improves performance and/or reduces dependence on extensive training. The final conclusion should be based on the measured trade-off between validation/test accuracy and computational cost. For the full 525-species problem, transfer learning is expected to be the most practical and scalable strategy.

\chapter{Appendix A: Example Code Snippet}
\begin{lstlisting}[language=Python, caption={Example Keras training call with class weighting}]
history = model.fit(
    train_flow,
    validation_data=valid_flow,
    epochs=20,
    callbacks=[early_stopping],
    class_weight=class_weights,
    verbose=1,
)
\end{lstlisting}

\chapter{Appendix B: Example Figure Placement}
\begin{figure}[H]
    \centering
    \includegraphics[width=0.8\textwidth]{example-image.png}
    \caption{Example model training curve or architecture figure. Replace with the actual image file.}
\end{figure}

\chapter{Appendix C: Example Table}
\begin{table}[H]
    \centering
    \caption{Example summary table for model comparison.}
    \begin{tabular}{lcccc}
        \toprule
        Model & Accuracy & Loss & Epochs & Training Time \\
        \midrule
        CNN A & -- & -- & -- & -- \\
        CNN B & -- & -- & -- & -- \\
        \bottomrule
    \end{tabular}
\end{table}

\end{document}
